\chapter{Results}
\label{ch:results}

This chapter reports and critically reviews the results of the project. It looks first at the accuracy of the two custom detection models, then at how the software behaved under automated testing, and finally at how the complete system performed when run end-to-end in real time. Results are reported honestly, including the weaknesses, because a fair account of what did \emph{not} work is as useful as the headline numbers.

\section{Evaluation Method}

All detection results were produced by a single, reproducible evaluation script that runs each trained model against its held-out test split, that is, data the model never saw during training. The standard object-detection measures are used: precision (of the things the model flagged, how many were correct), recall (of the things it should have found, how many it did), and mean Average Precision at an overlap threshold of 0.5 (mAP@0.5), together with the stricter mAP averaged over overlap thresholds from 0.5 to 0.95 (mAP@0.5:0.95). Software behaviour was checked with an automated test suite, and the whole system was exercised through scripted scenarios and live camera runs.

\section{Accessibility Detection Model}

The accessibility detector was tested on its 274-image held-out split. The overall and per-class results are given in Table~\ref{tab:acc_results}, and the confusion matrix in Figure~\ref{fig:cm_acc}.

\begin{table}[H]
\centering
\caption{Accessibility detection model results on the test split (YOLOv8s).}
\label{tab:acc_results}
\begin{tabular}{lcccc}
\toprule
\textbf{Class} & \textbf{Precision} & \textbf{Recall} & \textbf{mAP@0.5} & \textbf{mAP@0.5:0.95} \\
\midrule
door & 0.866 & 0.874 & 0.913 & 0.741 \\
door handle & 0.787 & 0.758 & 0.793 & 0.444 \\
up & 0.812 & 0.755 & 0.774 & 0.623 \\
updown & 1.000 & 0.000 & 0.000 & 0.000 \\
close & 0.888 & 0.827 & 0.895 & 0.664 \\
down & 0.833 & 0.777 & 0.812 & 0.598 \\
open & 0.821 & 0.851 & 0.895 & 0.654 \\
\midrule
\textbf{All classes} & \textbf{0.858} & \textbf{0.692} & \textbf{0.726} & \textbf{0.532} \\
\bottomrule
\end{tabular}
\source{Teamwork}
\end{table}

\begin{figure}[H]
\centering
\includegraphics[width=0.72\textwidth,keepaspectratio]{cm_accessibility.png}
\caption{Normalised confusion matrix for the accessibility detection model on the test split.}
\label{fig:cm_acc}
\source{Teamwork}
\end{figure}

The model does well on the classes that matter most for safety and interaction. Doors are detected strongly (mAP@0.5 = 0.913), and the lift ``close'' and ``open'' controls both reach 0.895. Door handles are harder (0.793), which is expected: a handle is a small object and small objects are always more difficult to localise precisely. Overall the detector reaches a mean precision of 0.858 and mAP@0.5 of 0.726.

One result deserves an honest critical note. The ``updown'' class collapses to zero recall: the model finds none of its test instances. This is not a mysterious failure but a data problem: that class is severely under-represented in the split, so the model never learned it well and there are too few examples to score against. This drags the overall recall (0.692) and mAP down and is the single clearest weakness in this model. The fix is straightforward in principle: collect and balance many more ``updown'' examples and retrain. It is worth stressing that in the running system these detections are not acted on blindly; the ByteTrack lifecycle layer requires an object to be seen consistently before it is trusted, which cushions the effect of occasional per-frame errors. For the assisted-interaction use case, this level of accuracy is usable, but the class imbalance is a genuine limitation to be addressed.

\section{Banknote Detection Model}

The Egyptian banknote detector was tested on a labelled held-out set of 266 images. Its results, in Table~\ref{tab:bank_results}, are excellent.

\begin{table}[H]
\centering
\caption{Banknote detection model results on the 266-image labelled test set (YOLOv8).}
\label{tab:bank_results}
\begin{tabular}{lcccc}
\toprule
\textbf{Class} & \textbf{Precision} & \textbf{Recall} & \textbf{mAP@0.5} & \textbf{mAP@0.5:0.95} \\
\midrule
1 EGP & 0.994 & 1.000 & 0.995 & 0.981 \\
1 EGP (coin) & 1.000 & 1.000 & 0.995 & 0.986 \\
5 EGP & 0.997 & 1.000 & 0.995 & 0.982 \\
10 EGP & 0.994 & 1.000 & 0.995 & 0.989 \\
new 10 EGP & 1.000 & 1.000 & 0.995 & 0.982 \\
20 EGP & 1.000 & 1.000 & 0.995 & 0.975 \\
new 20 EGP & 1.000 & 1.000 & 0.995 & 0.974 \\
50 piastres & 0.996 & 1.000 & 0.995 & 0.990 \\
50 EGP & 0.996 & 1.000 & 0.995 & 0.984 \\
100 EGP & 0.995 & 1.000 & 0.995 & 0.995 \\
200 EGP & 0.997 & 0.960 & 0.955 & 0.936 \\
\midrule
\textbf{All classes} & \textbf{0.997} & \textbf{0.996} & \textbf{0.991} & \textbf{0.979} \\
\bottomrule
\end{tabular}
\source{Teamwork}
\end{table}

The detector achieves near-perfect performance: an overall mAP@0.5 of 0.991, precision of 0.997, and recall of 0.996. Every denomination scores at least 0.955, and the confusion matrix (Figure~\ref{fig:cm_bank}) is almost purely diagonal, meaning notes are very rarely confused with one another, which is exactly what a user needs when the difference between two notes is real money. The only visible weakness is a small fraction of 200~EGP notes that go undetected (counted as background), reflected in that class's slightly lower recall of 0.960.

\begin{figure}[H]
\centering
\includegraphics[width=0.72\textwidth,keepaspectratio]{cm_banknote.png}
\caption{Normalised confusion matrix for the banknote detection model on the labelled test set.}
\label{fig:cm_bank}
\source{Teamwork}
\end{figure}

These strong numbers should be read with a little caution. The banknotes were photographed by the team, so the test set, while unseen during training, comes from a similar setting; performance on very different lighting, worn notes, or unusual angles may be lower. Even so, the near-perfect scores and the clean confusion matrix give real confidence that the custom dataset and detector are well matched to the task, and they clearly show the value of building a dataset specifically for the problem rather than relying on a generic model.

\section{Qualitative Results for the Other Modes}

The two detection models above can be scored with precise numbers because they have labelled test sets. The remaining modes, namely text reading, face recognition, and hand guidance, are harder to reduce to a single figure, so they were assessed qualitatively through their automated tests and live runs. Their results are reported here honestly, as observations rather than headline metrics.

\textbf{Text reading.} With the fine-tuned English model and the adaptive pre-processing pipeline, the reader handled clear, well-lit printed text reliably and spoke it in a sensible reading order. The stability buffer proved its worth in practice: by only reading text seen consistently across several frames, it stopped the system blurting out momentary misreads as the camera moved. The honest limitations are those flagged in the OCR-glasses literature: performance drops on small, distant, low-contrast, or steeply angled text, which is why the mode only activates when text is close and the user is still, and prompts the user to reposition if nothing readable is found.

\textbf{Face recognition.} The two-part matching rule, which requires both a strong similarity score and a clear margin over the second-best candidate, worked as intended in testing, reducing the chance of confidently naming the wrong person, which is the most damaging error this mode can make. Requiring a stable identity across several frames before speaking further reduced false announcements. The trade-off is that recognition is deliberately cautious: a briefly glimpsed or poorly lit face is more likely to be reported as ``unknown'' than guessed at, which for an assistive device is the safer failure.

\textbf{Hand guidance.} The interaction mode successfully translated the three-dimensional gap between the fingertip and a target object into directional vibration, moving through its guidance phases and firing a distinct confirmation pattern on success. Because it reuses the same object detections as navigation, its ceiling is partly set by the accessibility detector's accuracy on small objects such as handles, the weakest area of that model, so guidance is smoothest for larger, clearly detected targets.

\textbf{Audio and haptic feedback.} The priority queue behaved correctly under test: urgent obstacle warnings preempted routine announcements, and per-cue cooldowns prevented the same message repeating too quickly. Directional panning gave an intuitive left/right sense of where an object was. These are usability behaviours rather than accuracy metrics, but they are central to whether the system feels calm and trustworthy rather than chaotic.

\section{Software Testing}

Alongside the models, the application software was checked with an automated integration test suite run through \texttt{pytest}. These tests exercise the perception engines on sample inputs to confirm their end-to-end behaviour. As Table~\ref{tab:sw_tests} shows, all 22 automated tests passed, covering obstacle detection, OCR, interaction, banknote logic, and audio.

\begin{table}[H]
\centering
\caption{Automated integration (pytest) test results.}
\label{tab:sw_tests}
\begin{tabular}{lcc}
\toprule
\textbf{Test group} & \textbf{Cases} & \textbf{Result} \\
\midrule
Obstacle detection (runtime) & 7 & Pass \\
OCR (runtime) & 8 & Pass \\
Interaction detection (runtime) & 4 & Pass \\
Banknote detection (runtime) & 2 & Pass \\
Audio feedback (runtime) & 1 & Pass \\
\midrule
\textbf{Total} & \textbf{22} & \textbf{All pass} \\
\bottomrule
\end{tabular}
\source{Teamwork}
\end{table}

It is worth being clear about what these tests do and do not prove. They confirm that each engine behaves correctly on controlled inputs; they do not, on their own, prove that the whole system works together on a real camera. That is the purpose of the system-level testing described next.

\section{System Integration and Real-Time Performance}

To test Echora as a whole rather than as separate parts, two complementary methods were used. First, the state machine, the component that fuses every engine's output into one behaviour, was driven through scripted multi-frame scenarios that target the cross-engine interactions no single-engine test can reach: switch debouncing, minimum dwell time, the motion gate, and the danger override. Second, the complete system was run live on the camera in furnished indoor spaces while a user moved through the scene and triggered each mode.

The scripted scenarios all passed, as summarised in Table~\ref{tab:system_tests}. Most importantly, a danger-zone obstacle always forced an immediate return to navigation, whatever mode was active, confirming that the safety path really does take precedence.

\begin{table}[H]
\centering
\caption{System-level scenarios exercising the end-to-end state-machine behaviour.}
\label{tab:system_tests}
\begin{tabular}{cp{0.42\textwidth}p{0.24\textwidth}c}
\toprule
\textbf{\#} & \textbf{Scenario} & \textbf{Expected outcome} & \textbf{Result} \\
\midrule
1 & Idle scene, no triggers & Stay in navigation & Pass \\
2 & Text visible for a single frame & No switch (debounce) & Pass \\
3 & Stable text, user still, dwell elapsed & Switch to OCR & Pass \\
4 & Danger obstacle appears while in OCR & Override to navigation & Pass \\
5 & Text visible while user walks fast & Stay in navigation (motion gate) & Pass \\
6 & Stable, confident face detected & Switch to face mode & Pass \\
7 & Every mode change recorded & Transition logged & Pass \\
\bottomrule
\end{tabular}
\source{Teamwork}
\end{table}

During the live runs, the integrated system sustained real-time operation at roughly 20--21 frames per second while detecting and tracking several obstacles at once, estimating each object's distance and bearing, and sorting it into a safe, warning, or danger zone. Figure~\ref{fig:debug} shows two representative frames from the debug view. In the first, several chairs and a couch are tracked in the warning zone, each labelled with its distance and direction. In the second, a chair that has entered the danger corridor is flagged in red while a more distant couch stays green in the safe zone. The on-screen mode, frame counter, frame rate, and per-zone counts confirm that the whole pipeline, from capture through detection and tracking, depth zoning, and state evaluation to feedback, runs together in real time.

\begin{figure}[H]
\centering
\begin{minipage}[t]{0.49\textwidth}\centering
  \includegraphics[width=\linewidth,keepaspectratio]{debug_nav_warning.jpg}\\[2pt]
  {\footnotesize (a) Multiple obstacles in the warning zone}
\end{minipage}\hfill
\begin{minipage}[t]{0.49\textwidth}\centering
  \includegraphics[width=\linewidth,keepaspectratio]{debug_nav_danger.jpg}\\[2pt]
  {\footnotesize (b) Danger obstacle (red) vs.\ safe couch (green)}
\end{minipage}
\caption{Live end-to-end operation in navigation mode, from the real-time debug window. Each detection shows its class, distance, and bearing and is coloured by proximity zone; the display reports the active mode, frame rate ($\sim$20~FPS), and per-zone object counts.}
\label{fig:debug}
\source{Teamwork}
\end{figure}

\section{Discussion of Results}

Taken together, the results support the project's central claim: a unified, real-time, privacy-preserving assistive system can be built on ordinary hardware. The banknote model is genuinely strong; the accessibility model is usable but limited by one badly imbalanced class; the software passes its automated tests; and the full system runs in real time while respecting the safety-first rule.

Two honest limitations frame these results. First, the accuracy figures are the models' own test-set scores, not a measure of how helpful the system feels to a real BLV user; that requires user studies, which were beyond this prototype's scope. Second, the real-time runs confirm correct behaviour and frame rate, but a full account of latency, recognition accuracy in the wild, and safety-alert response time would need a dedicated hardware acceptance test. These are exactly the next steps set out in the following chapter.
